\section{问题定义}
\label{sec:setup}
\paragraph{部分可观测的社会动态。}
用 $i\in\{1,\ldots,N\}$ 表示智能体，$e$ 表示活动。时刻 $t$ 的全局状态为 $s_t=(h_{1t},\ldots,h_{Nt},g_t)$。其中，$h_{it}$ 包含兴趣、可用时间、剩余预算、已收信息和参与经历；环境状态 $g_t$ 包含通知、报名窗口、容量、报名结果、候补名单与日程。智能体获得局部观测 $o_{it}\sim O_i(\cdot\mid s_t)$，并形成意图 $u_{it}\sim\pi_\theta(\cdot\mid o_{it},h_{it})$。社会互动产生消息 $m_t\sim G_\theta(\cdot\mid s_t,u_t)$。环境把意图与消息结算为实际行动结果，再反馈到状态：
\begin{equation}
 a_t\sim K_\xi(\cdot\mid s_t,u_t,m_t),\qquad
 s_{t+1}\sim F_\xi(\cdot\mid s_t,a_t,m_t).
 \label{eq:loop}
\end{equation}
$\xi$ 表示环境规则。式~\eqref{eq:loop} 描述一个部分可观测的多智能体随机系统，并不要求额外声称它进行了奖励优化或策略学习。想报名不等于报名成功：容量、时间冲突和预算可能使意图无法兑现，取消与递补也会改变后续参与结果。

实现中的 $\theta=(a,b,d)$ 分别控制注意、条件参与和候补退出，在不同活动之间共享。个体状态和活动特征仍然不同，因此个体概率也不同。每个模拟世界中，同一批人参与相互影响的活动。群体轨迹 $X_{etc}$ 记录报名、候补和到场比例，$c\in\{1,2,3\}$ 表示三个通道。对候选 $m$ 与种子 $s$，缓存轨迹记为 $X_{msetc}$。每条轨迹全程使用同一个机制；平均参数或在时间步间切换机制，会定义另一个不同的模拟器。

\paragraph{历史信息与取证机会。}
历史数据包括观测轨迹，以及来自独立核查人群的二元阶段记录。取证池由另外一批人组成，每人只在一个标准化、机会状态已知的情境下贡献一条记录。候选概率为 $p_{mr}=P(Y_r=1\mid\theta_m,x_r)$，获取成本为 $k_r>0$。在本基准中，给定机制后各记录条件独立，记录是否可获取的元数据在查询前已知。我们不据此推断任意互动人群日志都相互独立。如果观察了组成记录，又把同一过程的汇总人数当成独立证据，就会重复计数；分离人群避免了这一特定依赖。

\paragraph{组合预测目标。}
活动配置记为 $\xi=(\ell,f,c,\eta)$，其中 $(\ell,f,c)$ 为候补、转发和冲突开关，$\eta$ 为其余属性。历史活动分别启用候补和转发；冲突规则存在于模拟器中，但所有历史配置均未启用。固定探针组合这些可访问规则，用于近似未来预测需求。取证策略 $\alpha$ 选择序列 $Q$，满足 $\sum_{r\in Q}k_r\leq B$，并更新机制混合分布。科学目标是在目标活动分布下获得较低的轨迹风险：
\begin{equation}
 \mathcal R_B(\alpha)=\E_{\xi,Y,Q\sim\alpha}\!\left[
 L\big(\widehat X(\xi\mid Y_Q),X^{\mathrm{true}}(\xi)\big)\right].
 \label{eq:risk}
\end{equation}
探针分歧是该风险的模型内部、可计算替代目标，并不保证是其无偏估计。缓存分为训练、探针、开发和测试库。开发阶段已包含成对和三重规则组合，测试采用不同活动实例，而不是整个开发阶段都未见过的组合。因此，本文研究的是超出经验训练活动的迁移，而非不受限制的组合分布外泛化。引入所有候选都没有的规则，则属于另一个模型失配实验。
